\documentclass{article}
\usepackage[UTF8]{ctex}
\usepackage{amsmath,amssymb,amsfonts}
\usepackage{enumitem}
\usepackage{geometry}
\geometry{a4paper, margin=1in}
% === 错题本排版优化（不参与题目标识） ===
\setlength{\parindent}{0pt}
\setlength{\parskip}{0.35em plus 0.10em minus 0.08em}
\setlength{\emergencystretch}{3em}
\setlist{itemsep=0.25em, topsep=0.40em, parsep=0pt}
\allowdisplaybreaks[2]
% === 错题本排版优化结束 ===
\begin{document}
\section*{太原理工大学2025年数学分析真题}
\section*{67．太原理工大学2025年数学分析试卷}

\subsection{题目1}
1．求极限 \(\lim _{x \rightarrow 0} \frac{\left(1+\sin ^{2} x\right)^{1011}-(\cos x)^{2026}}{\sqrt{1+2 x^{2}}-1}\) ．

\subsection{题目2}
2．求极限 \(\lim _{n \rightarrow \infty}\left[(n+2025)\left(\frac{1}{n^{2}+1}+\frac{1}{n^{2}+2^{2}}+\cdots+\frac{1}{n^{2}+n^{2}}\right)\right]\) 。

\subsection{题目3}
3．求不定积分 \(\int \frac{1+\sin x}{\sin x(1+\cos x)} \mathrm{d} x\) ．

\subsection{题目4}
4．求累次积分 \(\int_{0}^{1} x^{2} \mathrm{~d} x \int_{x}^{1} e^{-y^{2}} \mathrm{~d} y\) 。

\subsection{题目5}
5．计算曲面积分

\[
I=\iint_{\sum} 4 x z \mathrm{~d} y \mathrm{~d} z-2 y z \mathrm{~d} z \mathrm{~d} x+\left(1-z^{2}\right) \mathrm{d} x \mathrm{~d} y
\]

其中 \(S\) 是旋转曲面 \(z=e^{\sqrt{x^{2}+y^{2}}}\) 介于 \(z=1\) 和 \(z=e^{2}\) 的部分，取下侧．

\subsection{题目6}
6．求幂级数 \(\sum_{n=1}^{\infty}(-1)^{n-1} n^{2} x^{n}\) 的收敛域与和函数。

\subsection{题目7}
7．求 \(I(a)=\int_{0}^{+\infty} e^{-x^{2}} \cos (a x) \mathrm{d} x\) ．

\subsection{题目8}
8．讨论函数

\[
f(x, y)= \begin{cases}\frac{x^{3}-y^{3}}{x^{2}+y^{2}}, & x^{2}+y^{2} \neq 0 \\ 0, & x^{2}+y^{2}=0\end{cases}
\]

在原点 \((0,0)\) 处的连续性，偏导数的存在性，可微性．

\subsection{题目9}
9．\(f(x)\) 在 \([0,+\infty)\) 上非负连续，且

\[
f(x) \leq c+\int_{0}^{x} f(t) \mathrm{d} t, x \geq 0
\]

其中 \(c>0\) 为常数．证明：\(f(x) \leq c e^{x}, x \geq 0\) ．

\subsection{题目10}
10．证明：

\begin{itemize}
  \item[（1）] \(\frac{x}{1+x}<\ln (1+x)<x, x>0\) ．
  \item[（2）] \(\sum_{n=1}^{\infty} \frac{(-1)^{n}}{n} \cdot \frac{\ln \left(1+x^{n}\right)}{x^{n}}\) 在 \((0,+\infty)\) 上一致收敛．
\end{itemize}

\subsection{题目11}
11．\(f(x)\) 在 \([0,1]\) 上连续，\((0,1)\) 上可导，且 \(f(0)=f(1)=1\) ．证明：

\begin{itemize}
  \item[（1）] 存在 \(\eta \in(0,1)\) ，使得 \(e^{\eta}\left[f(\eta)+f^{\prime}(\eta)\right]=e-1\) ．
  \item[（2）] 存在 \(\xi \in(0,1)\) ，使得 \(e^{\eta-\xi}\left[f(\eta)+f^{\prime}(\eta)\right]=1\) ．
\end{itemize}

\subsection{题目12}
12．设 \(\left\{f_{n}(x)\right\}\) 在 \([a, b]\) 上一致收敛于 \(f(x)\) ，且对每一个正整数 \(n, f_{n}(x)\) 均在 \([a, b]\) 上连续．证明：

\begin{itemize}
  \item[（1）] \(f(x)\) 在 \([a, b]\) 上一致连续．
  \item[（2）] 对任意的 \(\varepsilon>0\) ，存在 \(\delta>0\) ，对任意的 \(x^{\prime}, x^{\prime \prime} \in[a, b]\) ，只要 \(\left|x^{\prime}-x^{\prime \prime}\right|<\delta\) ，对任意的正整数 \(n\) ，都有
\[
\left|f_{n}\left(x^{\prime}\right)-f_{n}\left(x^{\prime \prime}\right)\right|<\varepsilon .
\]
\end{itemize}
\end{document}
